% ch2-infrastructure.tex
\chapter{MLIR Provides the Infrastructure, ulm.lam Provides the Meaning}
\label{ch:infra}

\begin{answer}
MLIR~\cite{mlir} has an open, extensible type system. A dialect defines its own types,
attributes and operations, and states the rules that they must satisfy;
MLIR runs those rules at defined points and rejects IR that breaks them. MLIR
does not come with any particular type \emph{discipline}: not linear types,
not dependent types, not refinement types, not an effect system. Each of
those needs a semantics and an enforcement mechanism that someone has to
design and implement. In short: \emph{MLIR provides the infrastructure for
representing domain-specific types and enforcing rules about operations.
The meaning of those types, and the strength of the guarantees they provide,
comes from the dialect and its implementation.}
\end{answer}

The Field Guide describes the infrastructure side in \FG{3.1}{Dialects and
extensible operation sets} and \FG{3.2}{Traits and interfaces}, and the
type disciplines in \FGch{1}{Types as Static Resource Reasoning}. This
chapter shows, with ulm.lam, which of the two you are looking at in each
case.

% ---------------------------------------------------------------------------
\section{Who provides what}
\label{sec:infra-who}

\begin{table}[h]
\centering\small
\begin{tabularx}{\linewidth}{@{}P{2.4cm}LL@{}}
\toprule
\textbf{Concern} & \textbf{MLIR provides} & \textbf{ulm.lam implements} \\
\midrule
types & builtin \code{tensor} and \code{memref} with an \emph{encoding}
  slot; ODS (Operation Definition Specification, in TableGen) for new
  types and attributes &
  no new type; one attribute, \code{#ulm_lam.props<...>}, used as tensor
  encoding \\
operations & ODS: operands, results, type constraints, assembly format &
  \code{ulm_lam}: 7 ops (\code{mul}, \code{add}, \code{scale},
  \code{transpose}, \code{inverse}, \code{identity}, \code{assume});
  \code{ulm_lamk}: 17 kernel ops and \code{call} \\
verification & hooks (\code{verify()}, \code{verifyEncoding}), generated
  operand checks, \code{InferTypeOpInterface}'s result check; run after
  parsing and (by default) after each pass &
  the rules: inner dimensions, squareness, contradictory properties, DPS
  result types \\
inference & \code{InferTypeOpInterface} and its compatibility hook
  \code{isCompatibleReturnTypes} & the property rules (\code{mulProperties},
  \ldots) and what ``compatible'' means \\
effects & \code{MemoryEffectOpInterface}, the \code{Pure} trait; generic
  passes (CSE, DCE (dead code elimination), LICM (loop-invariant code
  motion)) that consult them & which ops are pure; which
  operands a kernel reads and writes \\
rewriting & pattern classes, the greedy driver, the dialect conversion
  framework & the normal form, kernel selection, the lowering to EmitC \\
memory & One-Shot Bufferize; \code{DstBufferizableOpInterfaceExternalModel}
  for DPS ops & DPS kernel ops and their bufferization models; the
  workspace and \code{info} passes \\
\bottomrule
\end{tabularx}
\caption{The division of labor between MLIR and ulm.lam.}
\label{tab:infra-who}
\end{table}

Read the table column by column. Everything on the left is generic and
knows nothing about matrices. Everything on the right is a decision of
ulm.lam. When this text says that ``the compiler knows that \(S\) is SPD'',
the knowledge is in the right column: MLIR stores the attribute, calls the
inference hook and runs the verifier, but the meaning of \code{spd} and the
rule that derives it are ulm.lam's C++.

% ---------------------------------------------------------------------------
\section{A property set as a tensor encoding}
\label{sec:infra-props}

\begin{steps}
\item[Input] a type written as \code{tensor<?x?xf64, #ulm_lam.props<spd>>}.
\item[Representation] a builtin \code{RankedTensorType} whose encoding is a \code{PropertiesAttr}, a 32-bit set of named facts.
\item[Rule] \code{verifyEncoding}, called by MLIR whenever such a type is built or parsed.
\item[Output] the type, or a diagnostic.
\item[Theory] a refinement of the base type by a predicate, \FG{1.7}{Refinement types}.
\end{steps}

The declaration (\srcref{include/ulm_lam/UlmLamAttrs.td}{37}{77}, cases
shortened):

\begin{lstlisting}[style=td]
def UlmLam_PropSquare     : I32BitEnumAttrCaseBit<"square",       0>;
def UlmLam_PropSymmetric  : I32BitEnumAttrCaseBit<"symmetric",    1>;
def UlmLam_PropSPSD       : I32BitEnumAttrCaseBit<"spsd",         2>;
def UlmLam_PropSPD        : I32BitEnumAttrCaseBit<"spd",          3>;
...
def UlmLam_PropFullRank   : I32BitEnumAttrCaseBit<"full_rank",   12>;
def UlmLam_PropLowerStored: I32BitEnumAttrCaseBit<"lower_stored",13>;
def UlmLam_PropUpperStored: I32BitEnumAttrCaseBit<"upper_stored",14>;
def UlmLam_PropDiagonal   : I32BitEnumAttrCaseGroup<"diagonal",
                              [UlmLam_PropLowerTri, UlmLam_PropUpperTri]>;
...
// Used as a builtin tensor encoding:  tensor<?x?xf64, #ulm_lam.props<spd>>
def UlmLam_PropsAttr : EnumAttr<UlmLam_Dialect, UlmLam_Properties, "props",
    [DeclareAttrInterfaceMethods<VerifiableTensorEncoding>]> {
  let assemblyFormat = "`<` $value `>`";
  ...
}
\end{lstlisting}

Two MLIR mechanisms carry it: \code{EnumAttr} over an \code{I32BitEnumAttr}
generates the attribute class, its parser and printer; the interface
\code{VerifiableTensorEncoding} makes MLIR call a verifier when a tensor
type with this encoding is created. The verifier is ulm.lam's
(\srcref{lib/UlmLamDialect.cpp}{145}{167}):

\begin{lstlisting}[style=cpp]
LogicalResult PropertiesAttr::verifyEncoding(ArrayRef<int64_t> shape, Type,
    function_ref<InFlightDiagnostic()> emitError) const
{
    P p = closure(getValue(), shape);
    if (shape.size() != 2)
        return emitError() << "#ulm_lam.props needs a rank-2 tensor";
    if (has(p, P::zero) && has(p, P::full_rank))
        return emitError() << "#ulm_lam.props: zero contradicts full rank";
    ...
    if (has(p, P::square) && !ShapedType::isDynamic(shape[0]) &&
        !ShapedType::isDynamic(shape[1]) && shape[0] != shape[1])
        return emitError() << "#ulm_lam.props: square property on a "
                           << shape[0] << "x" << shape[1] << " tensor";
    return success();
}
\end{lstlisting}

It first closes the set under implication (so \code{spd} brings
\code{square}), then checks the closure against itself and against the
static shape. Two rejected types (\code{outputs/verify.out}):

\begin{textout}
### v3-square-props.mlir
v3-square-props.mlir:1:53: error: #ulm_lam.props: square property on a 3x4 tensor
func.func @f(%a: tensor<3x4xf64, #ulm_lam.props<spd>>) {
### v4-zero-fullrank.mlir
v4-zero-fullrank.mlir:1:68: error: #ulm_lam.props: zero contradicts full rank
func.func @f(%a: tensor<?x?xf64, #ulm_lam.props<zero, non_singular>>) {
\end{textout}

\code{spd} on a \(3\times 4\) tensor is rejected because SPD implies square;
\code{zero} with \code{non_singular} because non-singular implies full rank.
Both are contradictions \emph{between the facts and the shape}, decided
without looking at any data. \code{tensor<4x4xf64, #ulm_lam.props<spd>>} is
accepted for every \(4\times 4\) matrix, SPD or not: a type cannot be wrong
about entries it never sees.

% ---------------------------------------------------------------------------
\section{Operation verifiers: structural rules}
\label{sec:infra-verify}

Each op class may define \code{verify()}. MLIR calls it after parsing and,
with the default pass-manager settings, after every pass. For
\code{ulm_lam.mul} (\srcref{lib/UlmLamDialect.cpp}{514}{526}):

\begin{lstlisting}[style=cpp]
LogicalResult MulOp::verify()
{
    if (getFactors().size() < 2)
        return emitOpError("needs at least two factors");
    for (unsigned i = 0; i + 1 < getFactors().size(); ++i) {
        Type a = getFactors()[i].getType(), b = getFactors()[i + 1].getType();
        if (!dimsCompatible(cols(a), rows(b)))
            return emitOpError("factor ") << i << " has " << cols(a)
                   << " columns but factor " << i + 1 << " has " << rows(b) << " rows";
    }
    return success();
}
\end{lstlisting}

\code{dimsCompatible(a, b)} is true if either is dynamic or both are equal.
Accepted and rejected (\code{outputs/verify.out}):

\begin{textout}
### v1-mul-static.mlir
v1-mul-static.mlir:2:8: error: 'ulm_lam.mul' op factor 0 has 4 columns but factor 1 has 5 rows
### v2-mul-dynamic.mlir
accepted
### v5-inverse-nonsquare.mlir
v5-inverse-nonsquare.mlir:2:8: error: 'ulm_lam.inverse' op operand is not square
### v9-gemm.mlir
v9-gemm.mlir:3:8: error: 'ulm_lamk.gemm' op shapes of op(A), op(B) and C do not match
### v10-gemm-dps.mlir
v10-gemm-dps.mlir:3:8: error: 'ulm_lamk.gemm' op result 0 must have the type of its init
\end{textout}

\code{v2} multiplies two \code{tensor<?x?xf64>}: the verifier cannot know
whether the inner sizes agree, so it accepts. \Cref{ch:shapes} shows where
dynamic sizes \emph{are} checked. \code{v10} violates a DPS rule: a kernel's
tensor result must have the type of the operand it overwrites
(\code{verifyDps}, \srcref{lib/UlmLamDialect.cpp}{826}{845}).

\begin{limit}
An operation verifier checks the IR, not the data. It sees types,
attributes and the structure of the program, and it runs at compile time.
It cannot tell whether the \(4\times 4\) array that a caller passes is SPD.
In ulm.lam that question is answered only by \code{potrf} at run time, and
only for the matrices that get factorized.
\end{limit}

% ---------------------------------------------------------------------------
\section{Result types must be justified}
\label{sec:infra-compat}

\begin{steps}
\item[Input] an op whose result type is written out, possibly with properties.
\item[Representation] \code{ulm_lam} ops implement \code{InferTypeOpInterface}.
\item[Rule] MLIR's generated verifier infers the result type from the operands and calls \code{isCompatibleReturnTypes(inferred, declared)}; ulm.lam defines compatible as ``same shape up to \code{?}, same element type, and every declared property is implied by the inferred ones''.
\item[Output] a declared type may be weaker than the inferred one, never stronger.
\item[Theory] subtyping by implication, \FG{1.7}{Refinement types}; checking against an annotation, \FG{1.10}{Type inference beyond Hindley–Milner: bidirectional and constraint-based}.
\end{steps}

The compatibility relation (\srcref{lib/UlmLamDialect.cpp}{359}{380}):

\begin{lstlisting}[style=cpp]
/// Shapes compatible, element types equal, properties of `actual` implied
/// by `inferred`.
bool ulmLamCompatible(TypeRange inferred, TypeRange actual)
{
    ...
        P have = closure(propertiesOf(it), it.getShape());
        if (!has(have, propertiesOf(at)))
            return false;
    ...
}
\end{lstlisting}

The test is the pair \code{v6} and \code{v7}: the same Gram product
\(A\T A\) of a \(4\times 3\) matrix, with the result declared \code{spd}.
Without \code{full_rank} on \(A\), only SPSD can be inferred, and the
verifier rejects the claim:

\begin{textout}
### v6-overclaim.mlir
v6-overclaim.mlir:3:8: error: 'ulm_lam.mul' op inferred type(s) 'tensor<3x3xf64, #ulm_lam.props<spsd>>' are incompatible with return type(s) of operation 'tensor<3x3xf64, #ulm_lam.props<spd>>'
v6-overclaim.mlir:3:8: error: 'ulm_lam.mul' op failed to infer returned types
### v7-overclaim-ok.mlir
accepted
\end{textout}

With \code{full_rank} on \(A\) (\code{v7}), SPD is inferable and the same
declaration is accepted. So inside a function, a property on a result type
is not an assumption: it is a claim that the verifier checks against the
inference rules. The only way to state more than inference derives is the
dedicated op \code{ulm_lam.assume}, which has no inference hook and is
accepted with any properties (\code{v8} in \code{outputs/verify.out}).
\Cref{sec:facts-declared} discusses it.

\begin{theory}[ (\FGs{1.7}, \FGs{1.10})]
In a refinement type system, a value of type \(\{v{:}B \mid p\}\) may be
used where \(\{v{:}B \mid q\}\) is expected if \(p \Rightarrow q\). Here
\(B\) is the tensor type, \(p\) the inferred property set and \(q\) the
declared one, and \(p \Rightarrow q\) is decided by closing \(p\) under the
implication rules and testing bit-set containment, not by an SMT
(satisfiability modulo theories) solver: with fifteen atoms and a fixed set
of Horn-like implications the check is a few bit operations. Comparing a
synthesized type against a written one is the \emph{checking} direction of
bidirectional typing.
\end{theory}

% ---------------------------------------------------------------------------
\section{Operation traits are not matrix properties}
\label{sec:infra-traits}

Both are called ``properties'' in casual speech, and both are written in
ODS. They live in different worlds.

\begin{center}\small
\begin{tabularx}{\linewidth}{@{}P{3.3cm}LL@{}}
\toprule
 & \textbf{Op traits and interfaces} & \textbf{Matrix properties} \\
\midrule
example & \code{Pure}, \code{Commutative}, \code{DestinationStyleOpInterface} & \code{spd}, \code{lower_tri}, \code{full_rank} \\
describes & an \emph{operation class} (all \code{ulm_lam.mul} ops) & a \emph{value} (this matrix) \\
where & the op's C++ class, via ODS & the value's type (encoding) \\
read by & generic MLIR passes: CSE, DCE, LICM, canonicalizer, bufferization & only ulm.lam's own rules and passes \\
MLIR understands it & yes: that is its purpose & no: an opaque attribute \\
\bottomrule
\end{tabularx}
\end{center}

\code{ulm_lam.add} has the trait \code{Commutative}
(\srcref{include/ulm_lam/UlmLamOps.td}{63}{71}); \code{ulm_lam.mul} does
not, because matrix products do not commute. The trait states a fact about
the \emph{operation} that a generic canonicalizer may use (to sort
operands). A matrix being SPD states a fact about the \emph{operand}, and no
upstream pass will ever read it.

\begin{limit}
MLIR operation traits are not mathematical matrix properties. The trait
\code{Pure} on \code{ulm_lam.inverse} says that the operation has no memory
effects and may be executed speculatively; it says nothing about whether the
operand is invertible. \Cref{sec:effects-speculation} shows what follows
from that difference.
\end{limit}

% ---------------------------------------------------------------------------
\section{What is present, and what is related by analogy}
\label{sec:infra-analogy}

\begin{table}[t]
\centering\small
\begin{tabularx}{\linewidth}{@{}P{3.2cm}LP{2.4cm}@{}}
\toprule
\textbf{Discipline (FG)} & \textbf{In ulm.lam} & \textbf{Status} \\
\midrule
refinement types (\FGs{1.7}) & property sets on tensor types; results
  checked by implication; parameters and \code{assume} trusted &
  \textbf{present}, with a finite predicate language \\
type inference (\FGs{1.10}) & synthesis of result types per op, one pass
  in program order & \textbf{present} \\
indexed / dependent types (\FGs{1.8}) & static shapes checked; run-time
  sizes named in the Julia frontend and in \code{ulm_lam.dims}; no
  type-level terms & \textbf{analogy} (\cref{ch:shapes}) \\
effect systems (\FGs{1.6}) & per-op memory effects, used by generic
  passes; no effects in function types, no effect polymorphism &
  \textbf{partly}: effect annotations, not an effect type system \\
linear / affine types (\FGs{1.2}, \FGs{1.3}) & none: SSA values may be
  used any number of times & \textbf{absent} (\cref{ch:values}) \\
uniqueness (\FGs{1.4}) & the in-place condition of bufferization, decided
  by an analysis & \textbf{analogy} (\cref{ch:storage}) \\
region types (\FGs{1.5}) & none; MLIR regions are nested code, the
  workspace is one arena per call & \textbf{absent}; analogy only \\
typestate (\FGs{1.9}) & ``factorized'' is tracked by a memo table in
  kernel selection, not in types & \textbf{absent} \\
\bottomrule
\end{tabularx}
\caption{Type disciplines of FG Chapter 1 and their status in ulm.lam.}
\label{tab:infra-status}
\end{table}

\Cref{tab:infra-status} is the short answer to ``which theoretical type
systems does ulm.lam have?''. Three entries need a remark.

\paragraph{Declared but not implemented} The ODS of \code{ulm_lam.assume}
has a unit attribute \code{checked}, documented as ``verified at run time
(e.g.\ symmetry test, or success of the Cholesky factorisation for spd)''
(\srcref{include/ulm_lam/UlmLamOps.td}{135}{154}). At commit
\texttt{\shortcommit} neither frontend sets it and no pass generates a
check from it; the only code that reads it is a normalization pattern that
leaves checked assumptions alone
(\srcref{lib/UlmLamPasses.cpp}{79}{99}). Likewise, the pass description in
\srcfile{include/ulm_lam/UlmLamPasses.td} lists kernel-selection strategies
\code{search} and \code{egraph}, a \code{cost-model} option and a pass
\code{convert-ulm-lamk-to-linalg}; that file is documentation (the passes
are registered by hand in C++), and the implemented pass rejects every
strategy but \code{gmc} (\srcref{lib/UlmLamPasses.cpp}{2295}{2305}).

\paragraph{Interfaces that nothing calls} \code{ulm_lam} and
\code{ulm_lamk} ops implement two interfaces of ulm.lam's own,
\code{PropertyInferenceOpInterface} and \code{KernelCostOpInterface}
(\srcref{include/ulm_lam/UlmLamInterfaces.td}{11}{48}). No pass at this
commit calls them: inference goes through MLIR's
\code{InferTypeOpInterface}, and kernel selection computes flop counts
inline (\cref{sec:orders-cost}). They are the extension points of a design,
not mechanisms of the current compiler.

\paragraph{Effects are partly present} MLIR's effect interface is a real
effect annotation per op, and generic passes use it soundly
(\cref{ch:values}). What the Field Guide calls an effect \emph{system},
effects in function types that compose by union and that a type checker
verifies, does not exist in MLIR or in ulm.lam.
